\documentclass[11pt]{article}

\usepackage[utf8]{inputenc}
\usepackage[T1]{fontenc}
\usepackage{lmodern}
\usepackage{amsmath}
\usepackage{amssymb}
\usepackage{booktabs}
\usepackage{graphicx}
\usepackage{xcolor}
\usepackage[hidelinks]{hyperref}
\usepackage{geometry}
\geometry{margin=1in}
\usepackage{natbib}
\bibliographystyle{plainnat}

\title{Does Reuse Survive Reliability? Lifecycle Optimization of Reuse-Based Servers for LLM-Inference Fleets}
\author{Anonymous Authors}
\date{}

\begin{document}
\maketitle

\begin{abstract}
Building servers from energy-efficient CPUs and reused DRAM and SSDs cuts the embodied carbon of cloud hardware, but such designs are evaluated at a single point in time, whereas operators decide over years whether to buy, reuse, cascade or retire hardware, and reused parts are already aged and fail more. We formulate a reliability-aware, multi-period lifecycle optimization for large-language-model (LLM) inference fleets that tracks hosts as cohorts by SKU, tier, region and age, precomputes age-dependent failures, repair drains and spare parts from Weibull hazards, couples reuse to a finite supply of harvested parts, allows cascading of previous-generation GPU hosts to a batch tier (in place or to a clean-grid region), and sizes latency-critical reserve with a Monte Carlo chance constraint. On a six-year, three-grid fleet whose parameters are assumed and swept (only the GPU failure rate is taken from published data, and GPU throughput, power and CXL slowdown are placeholders), reuse cut lifecycle carbon by 5.77\% on a clean grid but only 0.68\% and 0.10\% on mixed and dirty grids, because operational carbon was about 91\% of a GPU host's lifecycle; on mixed and dirty grids even the sign depends on the placeholder CXL slowdown. A reliability-aware plan retained 0.775\% of a 0.848\% fleet saving and degraded gracefully to zero as reused parts worsened, while a failure-blind plan retained 0.195\% and turned carbon-negative once reused parts failed about 2.1 times more than their age predicts. Guaranteeing a 0.1\% per-region-quarter shortfall probability under strongly correlated failures cost up to 0.82\% extra carbon, 83--85\% of it idle electricity of reserve hosts; the reuse saving rose slightly, to 0.812\%, because about a fifth of the reuse plan's reserve sat on hosts with lower embodied carbon and idle power. Cascading paid only near energy parity between old and new GPUs. We provide the measurement methodology to replace the placeholder parameters.
\end{abstract}

\section{Introduction}
Computing's carbon footprint now has two components of comparable concern: the operational carbon of the electricity that servers draw, and the embodied carbon emitted when their chips, memory, storage and chassis are manufactured \citep{gupta2022chasing,gupta2022act}. Machine learning sharpens both. Training and serving large models consume substantial energy \citep{strubell2019energy,patterson2021carbon,wu2022sustainableai}, end-to-end accounting of large language models (LLMs) must include the hardware they run on \citep{luccioni2023bloom,faiz2024llmcarbon}, and the carbon impact of generative-AI inference has been examined both today and in projection \citep{chien2023genai}. Because inference fleets are provisioned for years and refreshed continuously, how operators buy, reuse and retire hardware is itself a carbon lever.

One of the few levers demonstrated at cloud scale is \emph{component reuse}: building new server SKUs from energy-efficient CPUs together with DRAM and SSDs harvested from retired servers, with the reused memory attached over CXL \citep{wang2024greensku}. Reuse targets embodied carbon, which is large for storage and memory \citep{tannu2023ssd} and which architectural carbon models make explicit \citep{gupta2022act}; second-life hardware can be greener whenever the workload tolerates it \citep{switzer2023junkyard}.

Two properties of real fleets, however, are missing from such evaluations. First, \emph{decisions are sequential}. A reuse-based design is evaluated as a server at a point in time, under an average grid carbon intensity, while an operator decides every quarter, over many years, whether to buy new hosts, build reuse-based hosts, cascade older hosts to lighter service tiers, lease capacity or retire hardware---across regions whose grids differ and decarbonize at different rates. Second, \emph{reused parts are old}. DRAM and SSD failure rates vary with age and usage in the field \citep{schroeder2009dram,meza2015dram,schroeder2016flash}, and accelerator fleets fail frequently at scale \citep{llama3,kokolis2025reliability}. Every failure consumes a spare part, which carries embodied carbon, and drains node capacity for hours to days, which must be covered by headroom that itself carries embodied and idle carbon. Because a point-in-time evaluation neither tracks part age nor sizes headroom against concurrent failures, it cannot tell whether the embodied savings of reuse survive reliability. The gap widens for LLM inference: on 8-GPU hosts, operational carbon dwarfs embodied carbon, so the reach of an embodied-carbon lever depends on the grid.

We close this gap with a reliability-aware, multi-period lifecycle optimization for LLM-inference fleets. The model tracks the fleet as cohorts by SKU, service tier, region and age; it precomputes age-dependent expected failures, repair drains and spare parts from Weibull hazards so that the optimization stays linear; it couples reuse to a finite supply of harvested parts; and it allows cascading of previous-generation GPU hosts to a batch-inference tier, including shipping them to a clean-grid region. An outer Monte Carlo loop sizes per-region-quarter reserve so that a latency-critical service-level objective (SLO) holds with a target probability under independent or correlated failures. We evaluate a six-year, three-grid fleet; only the GPU failure rate is taken from published data, all other parameters are assumptions that we sweep, and GPU throughput, power and the slowdown of CXL-attached memory are placeholders pending measurements whose methodology we provide.

Our results answer the question in three parts. (1) The \emph{magnitude} of reuse savings is set by the grid: reuse cut fleet lifecycle carbon by 5.77\% on a clean grid but by 0.68\% on a mixed one and 0.10\% on a dirty one, because operational carbon was about 91\% of a new host's lifecycle carbon. (2) Reuse \emph{survives reliability if planned for}: a reliability-aware plan retained a saving of 0.775\% against 0.848\% when failures were ignored, and degraded gracefully to zero as reused parts became worse, whereas a failure-blind plan kept only 0.195\% and became carbon-negative once reused parts failed about twice as often as their age predicts. (3) Reuse \emph{persists under guaranteed SLOs}: meeting a 0.1\% per-region-quarter shortfall target under strongly correlated failures cost 0.79--0.82\% extra carbon, mostly idle electricity of reserve hosts, and the reuse saving grew slightly to 0.812\% because part of the reuse plan's reserve sat on hosts with lower embodied carbon and idle power. Cascading old GPUs paid only when their batch work per joule approached that of new GPUs.

Our contributions are: (i) a lifecycle model for reuse-based AI servers that jointly represents age-dependent component reliability, reused-part supply, cascading (in place and geographic) and regional grid decarbonization; (ii) break-even conditions for reused-part reliability and for cascading efficiency; (iii) chance-constrained SLO headroom validated by independent Monte Carlo replay, with the carbon price of the SLO; and (iv) design guidance for reliability-aware reuse in AI fleets, together with a measurement methodology to ground the remaining parameters.

\section{Background and Related Work}
\paragraph{Low-carbon server design and embodied-carbon accounting.}
Architectural carbon models quantify the embodied carbon of processors, memory and storage and expose its trade-off against operational efficiency \citep{gupta2022act,gupta2022chasing}; spatial-temporal embodied-carbon models refine how that carbon is attributed over a device's life \citep{zhang2024spatialembodied}, and datacenter-level frameworks co-design capacity with carbon-free energy supply \citep{acun2023carbonexplorer}. On the design side, reuse-based cloud servers combine efficient CPUs with reused DRAM over CXL and reused SSDs \citep{wang2024greensku}, SSD manufacturing has been identified as an embodied-carbon hot spot \citep{tannu2023ssd}, and discarded devices can serve as low-carbon compute \citep{switzer2023junkyard}. Circular-economy decisions---reuse, refurbishment, recycling and leasing---have been optimized with multi-objective fuzzy mathematical programs \citep{poonia2024circular}, building on fuzzy-set theory \citep{zadeh1965fuzzy}. These works evaluate designs or strategies without the age-dependent failure behaviour of reused parts and without the sequential procurement and retirement of a growing fleet. Our model embeds reuse-based SKUs in exactly that setting.

\paragraph{Field reliability of memory, storage and accelerators.}
Large field studies characterize DRAM errors and their dependence on age and platform \citep{schroeder2009dram,sridharan2012dram,sridharan2015memory,meza2015dram,zivanovic2019dram,dramicdcs2024}, SSD and flash failures \citep{meza2015flash,schroeder2016flash,ssdsystor2016,ssdnas2024}, and disk failure trends \citep{pinheiro2007disk}. Silent data corruption has emerged as a fleet-scale concern for processors \citep{hochschild2021cores,dixit2021sdc,sdctaco2024}, and large machine-learning clusters report frequent GPU and memory failures that interrupt jobs \citep{llama3,kokolis2025reliability}. These studies provide failure statistics but do not translate them into the carbon consequences of reusing aged parts. We calibrate age-dependent Weibull hazards from such statistics and propagate them into carbon through spares, repair drains and SLO headroom.

\paragraph{Carbon-aware scheduling and operation.}
A large body of work shifts \emph{when} and \emph{where} computation runs to follow grid carbon intensity: temporal shifting of deferrable work \citep{wiesner2021letswait}, carbon-aware capacity management in production datacenters \citep{radovanovic2023cics}, elastic scaling for carbon efficiency \citep{hanafy2024carbonscaler}, and virtual energy systems that expose carbon to applications \citep{souza2023ecovisor}; recent analysis also delineates the limits of temporal and spatial shifting \citep{sukprasert2024limitations}, and carbon-aware LLM inference adapts generation to grid conditions \citep{li2024sprout}. These approaches take the hardware fleet as given. We instead optimize the fleet itself over its lifecycle, which is complementary: geographic cascading moves \emph{hardware}, rather than work, to cleaner grids.

\paragraph{LLM serving, heterogeneous GPUs and tiered memory.}
LLM serving systems manage the key-value (KV) cache in paged memory \citep{kwon2023vllm}, split or disaggregate prefill and decode phases across machines \citep{patel2024splitwise,zhong2024distserve}, and exploit GPU heterogeneity for cost \citep{griggs2024melange}; power and energy management for LLM inference has been characterized and optimized \citep{patel2024power,stojkovic2024greener,stojkovic2025dynamollm}. These works optimize throughput, latency, cost or energy of a given fleet, not the lifecycle carbon of buying, reusing and cascading it. CXL enables memory pooling and tiering in cloud platforms \citep{li2023pond,berger2023cxlpooling,maruf2023tpp}, with measured characteristics on real devices \citep{sun2023demystifying} and debated benefits \citep{caseagainstcxl2023}; concurrent work tiers the KV cache across CXL memory \citep{rwcache2026}. On reliability, fault-injection studies show how bit errors propagate through DNN accelerators \citep{li2017dnnerrors}, and concurrent work examines bit-flip vulnerability of shared KV-cache blocks \citep{bitflipkv2026}. None of these considers aged, reused DIMMs as the home of the KV cache, or the carbon value of cascading old GPUs; we quantify the latter and provide a measurement methodology for the former. Our optimization is solved with a high-performance dual simplex solver \citep{huangfu2018highs}.

\section{Lifecycle Model}
\label{sec:model}
Figure~\ref{fig:overview} summarizes the model. It answers one question per quarter for every group of identical hosts: keep, retire, cascade, or complement with new or reuse-based hosts and leased capacity, so that demand is met at minimum lifecycle carbon.

\begin{figure}[!htbp]
  \centering
  \includegraphics[width=\linewidth]{figures/fig_lifecycle_model_overview.png}
  \caption{Overview of the reliability-aware lifecycle model. Server SKUs, per-age component reliability, regional grid carbon and demand, and reused-part pools feed a multi-period linear program over fleet cohorts that decides each quarter whether to buy, keep, retire, cascade (in place or to a clean-grid region) or lease capacity. An outer Monte Carlo loop adds latency-critical reserve until every region-quarter meets the target shortfall probability. Outputs are the reuse savings ladder, break-even analyses and the carbon price of a guaranteed SLO.}
  \label{fig:overview}
\end{figure}

\subsection{Fleet State, Decisions and Flow Conservation}
Time is discretized into quarters $t=0,\dots,T-1$ with $T=24$. The fleet is a set of cohorts indexed by SKU $s$, service tier $\tau\in\{\text{LC},\text{BATCH}\}$ (latency-critical and batch inference), region $r$ and node age $a$ in quarters within the node's current life; $n_{s\tau rat}\ge 0$ counts nodes. We model three 8-GPU inference SKUs: \texttt{new\_gpu}, a current-generation host with all-new components; \texttt{green\_gpu}, a reuse-based host in the style of \citet{wang2024greensku} with new GPUs, an energy-efficient CPU, and DRAM (attached over Compute Express Link, CXL) and SSDs that enter at age $a_0>0$ from a first life; and \texttt{old\_gpu}, a previous-generation host cascaded as a whole from its first life. Cohorts evolve as
\begin{equation}
n_{s\tau rat} = n_{s\tau r,a-1,t-1} - x^{\text{ret}}_{s\tau rat} - [\tau{=}\text{LC}]\sum_{r'} x^{\text{casc}}_{srr'at} + [\tau{=}\text{BATCH}]\sum_{r''} x^{\text{casc}}_{sr''rat} + [a{=}0]\,x^{\text{buy}}_{s\tau rt},
\label{eq:flow}
\end{equation}
with the initial fleet at $t=0$ and forced retirement at a maximum age. A cascade $x^{\text{casc}}_{srr'at}$ moves a node from the LC tier of region $r$ to the BATCH tier of region $r'$ and keeps its age; $r'\neq r$ is \emph{geographic cascading}, which incurs shipping carbon and cost. Building a reuse-based host consumes harvested DIMMs and SSDs from per-component pools, replenished by an external supply (retirements of a general-purpose fleet) and by first-life parts of hosts retired here, after a screening yield of 0.9.

\subsection{Reliability, Capacity and Carbon Accounting}
Each component class $c$ (GPU, CPU, DIMM, SSD) has a Weibull hazard
\begin{equation}
h_c(a)=\frac{k_c}{\eta_c}\left(\frac{a}{\eta_c}\right)^{k_c-1},\qquad H_c(a)=\left(\frac{a}{\eta_c}\right)^{k_c},
\label{eq:weibull}
\end{equation}
with part age $a$ in years (cohort ages in quarters are converted at 0.25 years per quarter), whose scale $\eta_c$ is calibrated so that the annual failure probability $1-e^{-(H_c(a^\star+1)-H_c(a^\star))}$ at a reference age $a^\star$ equals an annualized failure rate (AFR). A reused part starts at age $a_0$. A failed part is replaced by a new spare, so each component slot is a delayed renewal process whose first lifetime is conditioned on survival to $a_0$; we solve its renewal equation numerically. Because node age is a cohort index, the expected failures $\lambda_{sa}$, drained hours $d_{sa}$, availability loss $\ell_{sa}=d_{sa}/\Delta$ (with $\Delta$ hours per quarter), and spare and logistics carbon per node-quarter are \emph{precomputed}, and the optimization stays linear. A proportional-hazards multiplier $m$ scales the hazard of reused parts beyond what their age predicts.

Served work $w_{s\tau rat}$, measured in node-equivalents of a fully busy \texttt{new\_gpu} host, obeys
\begin{equation}
w_{s\tau rat}\le n_{s\tau rat}\,\theta_{s\tau}\,u\,(1-\beta_\tau\,\ell_{sa}),
\label{eq:cap}
\end{equation}
where $\theta_{s\tau}$ is relative throughput including the CXL multiplier of the reuse-based host, $u$ the utilization cap and $\beta_\tau$ an LC headroom multiplier ($\beta_{\text{BATCH}}=1$). LC demand $D_{rt}$ is per region; BATCH demand is global. Unserved demand is leased at a carbon premium or counted as shortfall. The objective is lifecycle carbon,
\begin{equation}
C=\sum\Big(e^{\text{emb}}_{sa}n + e^{\text{str}}_{sa}x^{\text{ret}}\Big)+\sum\Big(P^{\text{idle}}_{s}(1-\ell_{sa})\Delta\,n+\kappa_{s\tau}\,w\Big)\,\mathrm{PUE}\cdot\mathrm{CI}_{rt}+\sum\big(q^{\text{sp}}_{sa}+q^{\text{log}}_{sa}\big)n+C^{\text{lease}}+C^{\text{ship}},
\label{eq:obj}
\end{equation}
where embodied carbon $e^{\text{emb}}$ is amortized linearly over the design life and the unamortized remainder $e^{\text{str}}$ is charged on early retirement, operational carbon combines idle power with dynamic energy $\kappa_{s\tau}$ per unit of work and is scaled by the power usage effectiveness (PUE) of the facility, and the grid carbon intensity $\mathrm{CI}_{rt}$ declines per region. Embodied quantities follow component-level accounting in the spirit of architectural carbon models \citep{gupta2022act}. Reused parts carry zero embodied carbon (cut-off accounting) unless stated; we evaluate alternatives in Section~\ref{sec:robust}. Cost and shortfall serve as alternative objectives, and Pareto fronts are traced by $\varepsilon$-constraint. At fleet scale we solve the linear-programming relaxation with HiGHS \citep{huangfu2018highs}.

\subsection{Chance-Constrained SLO Headroom}
Equation~\eqref{eq:cap} protects the LC tier against \emph{expected} drains. An SLO, however, concerns the instantaneous number of drained nodes. With failures arriving at rate $\lambda$ and each repair lasting $d$ hours, the number of drained nodes of a cohort of $N$ nodes is the occupancy of an M/G/$\infty$ queue: Poisson with mean $N\,d_{sa}/\Delta$ for independent failures, or negative binomial with dispersion $k$ for correlated failures (variance $\mu+\mu^2/k$). We add a reserve $R_{rt}$ on \emph{installed} LC capacity,
\begin{equation}
\sum_{s,a}\theta_{s,\text{LC}}\,u\,(1-\beta_{\text{LC}}\ell_{sa})\,n_{s,\text{LC},rat}+\text{lease}_{rt}\ \ge\ D_{rt}+R_{rt},
\label{eq:reserve}
\end{equation}
and iterate: solve; sample instantaneous capacity per region-quarter under the planned fleet; set $R_{rt}\leftarrow R_{rt}+1.05\,\max\{0,\,D_{rt}-Q_\varepsilon(\text{capacity}_{rt})\}$, where $Q_\varepsilon$ is the $\varepsilon$-quantile; re-solve until every gap is at most 0.05\% of demand. Reserves grow monotonically, so the loop terminates; reserve nodes are idle and therefore pay idle power and embodied carbon, which is the carbon price of the SLO. Final plans are validated by an independent replay with 50{,}000 samples.

\section{Evaluation Methodology}
\label{sec:method}
\subsection{Scenarios and Parameters}
\emph{Node level:} one host per SKU over a six-year life at 60\% utilization on three grids with base carbon intensity 40, 250 and 650\,gCO$_2$e/kWh (clean, mixed, dirty) declining 2\%, 5\% and 3\% per year, and PUE 1.2. \emph{Reference fleet:} 24 quarters and three regions; LC demand starts at 180, 450 and 330 node-equivalents (clean, mixed, dirty) and grows 25\%, 25\% and 30\% per year; global BATCH demand starts at 150 and grows 35\% per year. The initial fleet holds about 1{,}750 hosts: 150/500/400 previous-generation hosts and 100/350/250 current-generation hosts per region. Reused-part pools start with 3{,}200 DIMMs and 800 SSDs and receive 2{,}400 DIMMs and 600 SSDs per quarter.

Per-unit embodied carbon is 150\,kgCO$_2$e for a current-generation GPU, 120 for a previous-generation GPU, 30 for a CPU, 20 for a 64\,GB DIMM and 614 for a 3.84\,TB SSD; \texttt{new\_gpu} draws 2{,}500\,W idle and 10{,}200\,W peak. Weibull shapes and AFRs are 1.0/0.0868 (GPU), 1.0/0.002 (CPU), 1.6/0.01 (DIMM) and 1.3/0.01 (SSD), with drains of 72, 24, 8 and 2 hours per failure. The GPU AFR is derived from a published interruption report of a 54-day pretraining run on about 16{,}384 GPUs with 148 faulty-GPU and 72 GPU-memory interruptions \citep{llama3}. All other parameters are engineering assumptions with ranges that we sweep; Table~\ref{tab:params} in Appendix~\ref{app:extra} lists the remaining model settings. \textbf{GPU throughput, power and the CXL slowdown are placeholders} (relative throughput of \texttt{old\_gpu} 0.30 LC and 0.45 BATCH; CXL multiplier 0.95 LC and 0.98 BATCH); results that depend on them are model-based. No proprietary fleet data is used.

\subsection{Experiment Ladder, Metrics and Solver}
We compare plans on the same scenario. E0: no reuse (only new hardware), with failures ignored (E0$_\text{det}$) or counted. E1: reuse allowed, failures ignored. E2a: E1's purchase and retirement decisions replayed with failures. E2b: reuse with a reliability-aware plan. E3/E3g: E2b plus in-place / geographic cascading. Savings are $X=1-C(\text{E1})/C(\text{E0}_\text{det})$, $Y_\text{naive}=1-C(\text{E2a})/C(\text{E0})$, $Y=1-C(\text{E2b})/C(\text{E0})$, and $Z$ ($Z_g$) is the additional saving of E3 (E3g) over E2b in percentage points.

The mixed-integer formulation of a 36-node pilot hit a 300\,s limit without proving optimality, while its LP relaxation solved in under 1\,s, matched the integer $X$ and $Y$ within about 0.05 percentage points, and gave identical results at 1$\times$ and 50$\times$ fleet scale; all fleet results therefore use the relaxation, and a full ladder solves in about 4\,s. A suite of 40 tests, including a hand-calculated single-SKU case, passes.

We also prepared, but have not yet executed, three GPU measurements to replace the placeholders: a serving benchmark on a paged-KV-cache serving engine \citep{kwon2023vllm} that measures LC goodput under time-to-first-token and time-per-output-token SLOs and BATCH throughput with board power (tokens per joule) per GPU generation; a KV-cache offload benchmark with the cache in local versus remote non-uniform memory access (NUMA) host memory as a proxy for CXL-attached reused DRAM; and KV-cache fault injection at controlled bit-error rates under no protection, single-error-correcting double-error-detecting (SECDED) ECC, and block checksums with recomputation.

\section{Results}
\label{sec:results}
\subsection{Operational Carbon Dominates GPU Hosts}
Table~\ref{tab:node} gives six-year node-level carbon on the mixed grid. Operational carbon was about 91\% of a new host's lifecycle carbon. Reuse cut the host's embodied carbon by about two thirds, from 8{,}312 to 2{,}860\,kgCO$_2$e, but changed its total by only about 6\%. Per unit of work, the reuse-based host saved 18.23\% (LC) and 20.73\% (BATCH) on the clean grid, 1.26\% and 4.28\% on the mixed grid, and $-2.04\%$ and 1.08\% on the dirty grid with failures ignored; counting failures lowered these by 0.1 to 2.2 percentage points. The mixed-grid result was dominated by the CXL throughput multiplier, whose range of 0.85--1.0 moved the saving from $-10.72\%$ to 5.88\%, a swing of 16.61 percentage points---three times the next parameter, the reuse host's peak-power delta (5.65), followed by SSD embodied carbon (5.41) and grid carbon intensity (4.20) (Figure~\ref{fig:tornado}).

\begin{table}[!htbp]
  \centering
  \caption{Six-year lifecycle carbon per node (kgCO$_2$e), mixed grid, LC tier, failures counted.}
  \label{tab:node}
  \begin{tabular}{lrrrrr}
    \toprule
    SKU & Embodied & Operational & Spares + logistics & Total $\downarrow$ & Failures/node \\
    \midrule
    \texttt{new\_gpu} & 8{,}312 & 96{,}662 & 1{,}291 & 106{,}265 & 10.62 \\
    \texttt{green\_gpu} & 2{,}860 & 95{,}511 & 1{,}575 & 99{,}945 & 15.18 \\
    \texttt{old\_gpu} & 0 & 62{,}344 & 1{,}896 & 64{,}240 & 18.42 \\
    \bottomrule
  \end{tabular}
\end{table}

\begin{figure}[!htbp]
  \centering
  \includegraphics[width=0.68\linewidth]{figures/fig_node_sensitivity_tornado.png}
  \caption{One-at-a-time sensitivity of the node-level reuse saving (mixed grid, latency-critical tier) to each parameter's range. Bars show the saving with the parameter at the low (blue) and high (orange) ends of its range. The CXL throughput multiplier has by far the largest swing; several parameters can flip the sign of the small base saving.}
  \label{fig:tornado}
\end{figure}

\subsection{Does Reuse Survive Reliability?}
Table~\ref{tab:ladder} and Figure~\ref{fig:ladder} report the fleet ladder. Ignoring failures, reuse saved $X=0.848\%$. A reliability-aware plan kept $Y=0.775\%$, about 91\% of it. A failure-blind plan kept only $Y_\text{naive}=0.195\%$: it purchased the same reuse-heavy fleet as E1, then had to lease 13{,}104\,tCO$_2$e of capacity to cover drained nodes, against 530\,tCO$_2$e for the reliability-aware plan. Split by grid, $X$ was 5.766\%, 0.675\% and 0.103\% and $Y$ was 5.028\%, 0.568\% and 0.087\% on clean, mixed and dirty single-grid fleets, while the failure-blind plan saved 4.584\% and 0.039\% on clean and mixed grids and lost 0.443\% on the dirty grid.

\begin{table}[!htbp]
  \centering
  \caption{Reference fleet: six-year lifecycle carbon by plan (tCO$_2$e).}
  \label{tab:ladder}
  \begin{tabular}{lrrrrr}
    \toprule
    Plan & Total $\downarrow$ & Embodied & Operational & Spares + log. & Leased \\
    \midrule
    E0$_\text{det}$: no reuse, failures ignored & 484{,}413 & 24{,}134 & 460{,}279 & 0 & 0 \\
    E0: no reuse, failures counted & 490{,}607 & 24{,}606 & 462{,}068 & 3{,}933 & 0 \\
    E1: reuse, failures ignored & 480{,}305 & 18{,}871 & 460{,}937 & 0 & 497 \\
    E2a: failure-blind plan, replayed & 489{,}652 & 18{,}871 & 453{,}457 & 4{,}220 & 13{,}104 \\
    E2b: reliability-aware reuse & 486{,}804 & 19{,}716 & 462{,}295 & 4{,}264 & 530 \\
    E3: E2b + cascading & 486{,}780 & 19{,}715 & 462{,}574 & 4{,}301 & 191 \\
    \bottomrule
  \end{tabular}
\end{table}

\begin{figure}[!htbp]
  \centering
  \includegraphics[width=\linewidth]{figures/fig_fleet_carbon_ladder_by_grid.png}
  \caption{Left: change in six-year lifecycle carbon relative to the no-reuse fleet, by carbon component, for a failure-blind reuse plan, a reliability-aware reuse plan, and the latter with in-place or geographic cascading; dots mark the net change. Right: reuse savings $X$ (failures ignored), failure-blind, $Y$ (reliability-aware) and $Y$ plus cascading for single-grid fleets on clean, mixed and dirty grids. GPU throughput, power and CXL slowdown are placeholder values.}
  \label{fig:ladder}
\end{figure}

Figure~\ref{fig:reliab} sweeps the reused-part hazard multiplier $m$. The reliability-aware plan bought fewer reuse-based hosts as $m$ grew, with $Y$ falling from 0.863\% ($m=0.5$) to 0.621\% ($m=2$), 0.209\% ($m=5$) and 0 from $m=8$ onwards. Because the planner can always fall back to the no-reuse fleet, $Y\ge 0$ holds by construction; the informative result is how fast $Y$ decays, and that a plan made without reliability information has no such floor. The failure-blind plan crossed zero near $m\approx2.1$ and reached $-2.251\%$ at $m=20$. At node level, the aged DIMMs nearly doubled DIMM failures (9.67 against 5.43 per host over six years), but each DIMM swap drains 8\,h whereas GPU failures, common to both SKUs, drain 72\,h; GPU failures therefore accounted for most drained hours, and total drained hours rose only from 359 to 394 per host, which is why reliability eroded savings mainly through planning rather than through raw failure counts.

\begin{figure}[!htbp]
  \centering
  \includegraphics[width=0.6\linewidth]{figures/fig_reliability_breakeven_reused_parts.png}
  \caption{Fleet carbon saving of reuse versus a multiplier on reused-part failure rates beyond the age effect. The reliability-aware plan degrades gracefully and stops reusing at a multiplier of about 8, never doing worse than no reuse; a plan that ignores failures crosses zero near 2 and becomes increasingly carbon-negative.}
  \label{fig:reliab}
\end{figure}

\subsection{When Does Cascading Pay?}
Cascading added essentially nothing at the placeholder efficiency ($Z=0.005$ percentage points; old/new BATCH work per joule 0.70). Figure~\ref{fig:cascade} sweeps that ratio. $Z$ stayed at or below 0.027 points up to a ratio of 0.935, then rose to 0.109 (ratio 1.168) and 0.346 (ratio 1.869). Above parity, geographic cascading to the clean region's BATCH tier added roughly 50--70\% on top ($Z_g=0.184$ and 0.521 at ratios 1.168 and 1.869). Because cascading avoids only the embodied carbon of new hosts, which is a small share of a GPU host's lifecycle, it cannot compensate for an operational-efficiency gap.

\begin{figure}[!htbp]
  \centering
  \includegraphics[width=0.6\linewidth]{figures/fig_cascade_breakeven_energy_efficiency.png}
  \caption{Additional carbon saving from cascading previous-generation GPU hosts to the batch tier versus their batch work per joule relative to current-generation hosts. Cascading pays only near energy parity; shipping cascaded hosts to the clean-grid region adds further savings once cascading pays.}
  \label{fig:cascade}
\end{figure}

\subsection{Guaranteeing the SLO}
Replaying the expected-value plans (headroom multiplier 3) against sampled drains showed that the LC tier met demand on average but not everywhere: the worst region-quarter had a 2.7\% probability of LC shortfall without reuse and 14.5\% with reliability-aware reuse under independent failures, rising to 15.5\% under correlated failures; under strong correlation ($k=0.5$), 44--48 of 72 region-quarters exceeded 1\%. Table~\ref{tab:slo} and Figure~\ref{fig:slo} report the chance-constrained plans. In independent validation the worst region-quarter shortfall probability was 0.96--1.07\% for a 1\% target and 0.10--0.13\% for a 0.1\% target across reuse and no-reuse plans, i.e. at the target within sampling error (the standard error at 0.1\% with 50{,}000 samples is about 0.014 percentage points) and the loop's stopping tolerance of 0.05\% of demand; the loop converged in 4 to 15 iterations. The SLO cost 0.22--0.82\% extra carbon, rising with strictness and correlation; strongly correlated failures required about 2.6$\times$ the reserve of independent ones at the 0.1\% target. In every setting the reuse plan paid a slightly lower price than the no-reuse plan, so the reuse saving \emph{grew} under the guarantee: from 0.775\% to 0.812\%, and from 0.480\% to 0.500\% when reused parts failed three times more than their age predicts. (The expected-value baseline here uses $\beta_{\text{LC}}=1$, which is why $Y$ at $m=3$ is 0.480\% rather than the 0.475\% of the $\beta_{\text{LC}}=3$ sweep.)

To identify the mechanism, Table~\ref{tab:reserve} decomposes the strictest setting ($\varepsilon=0.001$, $k=0.5$, $m=1$). The SLO's added carbon was dominated by idle electricity of reserve hosts (83\% without reuse, 85\% with reuse), not by embodied carbon. The reuse plan held 21.7\% of its added reserve capacity on reuse-based hosts---less than their 30\% share of the fleet, so the planner did not preferentially use them as reserve---and its reserve emitted 214\,tCO$_2$e (5.3\%) less, split between 111\,tCO$_2$e less embodied and 115\,tCO$_2$e less operational carbon, reflecting the lower embodied carbon and slightly lower idle power of reuse-based hosts. The effect is small (0.037 percentage points) and should be read as ``reuse does not lose under a guarantee'' rather than as a strong gain.

\begin{table}[!htbp]
  \centering
  \caption{Composition of the added SLO reserve ($\varepsilon=0.001$, $k=0.5$, $m=1$): chance-constrained minus expected-value plan. Capacity in node-equivalents summed over region-quarters; carbon in tCO$_2$e.}
  \label{tab:reserve}
  \small
  \begin{tabular}{lrrrrrr}
    \toprule
    Plan & New & Reuse-based & Prev.-gen. & Added carbon $\downarrow$ & Operational & Embodied \\
    \midrule
    no reuse & 1{,}320.3 & 0 & 0 & 4{,}020 & 3{,}356 & 575 \\
    reuse & 965.6 & 276.4 & 30.6 & 3{,}806 & 3{,}241 & 465 \\
    \bottomrule
  \end{tabular}
\end{table}

\begin{table}[!htbp]
  \centering
  \caption{Chance-constrained LC headroom: reuse saving $Y$ under a guaranteed SLO and the carbon price of the SLO, both in \%. $m$ is the reused-part failure multiplier and $\varepsilon$ the per-region-quarter shortfall target.}
  \label{tab:slo}
  \begin{tabular}{rlrrrrr}
    \toprule
    $m$ & Correlation & $\varepsilon$ & $Y$ guaranteed $\uparrow$ & $Y$ expected & Price (reuse) & Price (no reuse) \\
    \midrule
    1 & independent & 0.01 & 0.786 & 0.775 & 0.217 & 0.229 \\
    1 & independent & 0.001 & 0.788 & 0.775 & 0.301 & 0.314 \\
    1 & $k=2$ & 0.01 & 0.786 & 0.775 & 0.292 & 0.303 \\
    1 & $k=2$ & 0.001 & 0.792 & 0.775 & 0.427 & 0.445 \\
    1 & $k=0.5$ & 0.01 & 0.795 & 0.775 & 0.477 & 0.497 \\
    1 & $k=0.5$ & 0.001 & 0.812 & 0.775 & 0.785 & 0.823 \\
    3 & independent & 0.01 & 0.485 & 0.480 & 0.223 & 0.229 \\
    3 & independent & 0.001 & 0.484 & 0.480 & 0.310 & 0.314 \\
    3 & $k=2$ & 0.01 & 0.484 & 0.480 & 0.299 & 0.303 \\
    3 & $k=2$ & 0.001 & 0.490 & 0.480 & 0.435 & 0.445 \\
    3 & $k=0.5$ & 0.01 & 0.488 & 0.480 & 0.489 & 0.497 \\
    3 & $k=0.5$ & 0.001 & 0.499 & 0.480 & 0.803 & 0.823 \\
    \bottomrule
  \end{tabular}
\end{table}

\begin{figure}[!htbp]
  \centering
  \includegraphics[width=0.9\linewidth]{figures/fig_chance_constrained_slo.png}
  \caption{Left: reuse saving under a guaranteed latency-critical SLO (per region-quarter shortfall probability of at most 1\% or 0.1\%) for independent and increasingly correlated failures, with reused parts failing as their age predicts (1$\times$) or three times more (3$\times$). Right: carbon price of the SLO for the reuse plan relative to its expected-value plan.}
  \label{fig:slo}
\end{figure}

\subsection{Robustness: Accounting, Uncertainty, Headroom and Cost}
\label{sec:robust}
\emph{Accounting.} The convention for reused parts changed magnitudes but not rankings (Table~\ref{tab:robust}, Appendix~\ref{app:extra}): $Y$ was 0.775\% with cut-off and amortized accounting, 0.601\% when reused parts carried their unamortized share, 1.423\% and 1.067\% with upfront accounting, and 0 when reused parts were charged as new, in which case the failure-blind plan lost 0.506--0.557\%. \emph{Fuzzy uncertainty.} With triangular factors on failure rates $\times[0.5,1,3]$, embodied carbon $\times[0.5,1,2]$ and the CXL multiplier $[0.85,0.95,1.0]$, propagated by $\alpha$-cuts \citep{zadeh1965fuzzy} as in multi-objective fuzzy circular-economy models \citep{poonia2024circular}, $Y$ lay in $[0.000,2.638]\%$ at $\alpha=0$ and $[0.266,1.489]\%$ at $\alpha=0.5$. The lower bound of zero is guaranteed by the no-reuse fallback; the width of the interval, driven mainly by embodied carbon, is the informative result: at full support the saving is anywhere between nothing and about three times the nominal value. \emph{Utilization and headroom.} $Y$ was higher at a 0.5 utilization cap (0.783--0.806\%) than at 0.8 (0.775--0.777\%), consistent with embodied carbon mattering more on lightly loaded capacity; the share of node-quarters on reuse-based hosts fell from 0.303 to 0.280 as the LC headroom multiplier rose from 1 to 10. \emph{Cost.} The cost-minimal plan emitted 499.5\,ktCO$_2$e at USD 1.437 billion and the carbon-minimal plan 486.8\,ktCO$_2$e at USD 1.719 billion: 2.6\% more carbon for 16\% less cost, with most of the carbon gain available within about 5\% extra cost (Figure~\ref{fig:pareto}, Appendix~\ref{app:extra}).

\section{Discussion and Limitations}
\label{sec:discussion}
\subsection{Design Implications}
Three implications follow for reliability-aware reuse in AI fleets. First, place reuse where it pays: reuse-based hosts belong in clean-grid regions and in lightly loaded capacity, where embodied carbon weighs most, while steadily loaded capacity on carbon-intensive grids is best served by the most energy-efficient hosts. Reserve capacity is a special case: its carbon is mostly idle electricity, so powering reserve hosts down when not needed matters more than which SKU holds it. Second, plan for reliability rather than ignoring it: the same reuse decision is beneficial under a reliability-aware plan and harmful under a failure-blind one, so lifecycle planners must carry per-age failure, drain and spare models. Third, treat cascading old GPUs as a cost lever unless their batch energy efficiency approaches parity with new GPUs; if it does, ship them to clean grids.

\subsection{Limitations}
Our conclusions are bounded by the modelling setting. GPU throughput, power and the CXL slowdown are placeholders; the cascading verdict and the sign of reuse savings on mixed grids depend on them, which is why we provide the measurement methodology of Section~\ref{sec:method}. Except for the GPU failure rate, embodied-carbon and failure parameters are assumptions whose ranges we sweep. Reliability enters the LP as expected values, with tail behaviour handled by the outer Monte Carlo loop rather than inside the optimization, and fleet results use the LP relaxation. The study uses no proprietary fleet data.

\section{Conclusion}
\label{sec:conclusion}
We asked whether the carbon savings of server component reuse survive the reliability of aged, reused parts in LLM-inference fleets. A reliability-aware lifecycle optimization answers in three parts: the grid sets the magnitude of reuse savings; reuse survives reliability when planning accounts for it and fails when it does not; and reuse persists under guaranteed SLOs, whose carbon price is mostly idle electricity of reserve hosts. Executing the prepared GPU measurements is the next step to replace the remaining placeholders.

\bibliography{refs}

\clearpage
\appendix
\section{Additional Settings and Results}
\label{app:extra}

\begin{table}[!htbp]
  \centering
  \caption{Remaining model settings (assumptions unless stated).}
  \label{tab:params}
  \small
  \begin{tabular}{ll}
    \toprule
    Setting & Value \\
    \midrule
    Utilization cap $u$; LC headroom multiplier $\beta_{\text{LC}}$ & 0.8; 3 (1 in the chance-constrained study) \\
    \texttt{old\_gpu} power (idle / peak); previous-gen.\ GPU Weibull shape & 1{,}800 / 6{,}500\,W; 1.3 \\
    \texttt{green\_gpu} power vs \texttt{new\_gpu} (idle / peak) & $-50$ / $-100$\,W \\
    Platform embodied carbon (\texttt{new}, \texttt{old} / \texttt{green}) & 1{,}500 / 1{,}600\,kgCO$_2$e \\
    Design life; maximum age (\texttt{new}, \texttt{green} / \texttt{old}) & 6 years; 32 / 16 quarters \\
    Leased capacity carbon; shipping between regions & 1.3$\times$ new host per unit work; 300\,kgCO$_2$e/node \\
    Capital cost (\texttt{new} / \texttt{green} / \texttt{old}) & USD 300{,}000 / 280{,}000 / 15{,}000 \\
    Electricity; lease; cascade; shipping & USD 0.08/kWh; 60{,}000 per node-eq.-quarter; 3{,}000; 5{,}000 \\
    Repair cost per failure (GPU / prev.\ GPU / CPU / DIMM / SSD) & USD 2{,}500 / 4{,}000 / 800 / 300 / 400 \\
    Failure correlation $k$ (negative binomial) & 2 and 0.5, illustrative, not fitted \\
    \bottomrule
  \end{tabular}
\end{table}

\begin{table}[!htbp]
  \centering
  \caption{Reuse savings (\%) under embodied-carbon accounting conventions for reused parts.}
  \label{tab:robust}
  \begin{tabular}{llrrr}
    \toprule
    Reused-part charge & Amortization & $X$ & $Y_\text{naive}$ & $Y$ \\
    \midrule
    cut-off (zero) & amortized & 0.848 & 0.195 & 0.775 \\
    cut-off (zero) & upfront & 1.544 & 0.918 & 1.423 \\
    residual share & amortized & 0.667 & 0.029 & 0.601 \\
    residual share & upfront & 1.181 & 0.562 & 1.067 \\
    charged as new & amortized & 0.000 & $-0.557$ & 0.000 \\
    charged as new & upfront & 0.000 & $-0.506$ & 0.000 \\
    \bottomrule
  \end{tabular}
\end{table}

\begin{figure}[!htbp]
  \centering
  \includegraphics[width=0.5\linewidth]{figures/fig_carbon_cost_pareto.png}
  \caption{Carbon--cost $\varepsilon$-constraint front of the reference fleet over six years.}
  \label{fig:pareto}
\end{figure}

\end{document}
